% open-logic.tex
% compiles to produce complete text using open-logic-dev style

\documentclass[../include/open-logic-part]{subfiles}

\begin{document}

\clearpage

\begin{editorial}
ایہ فائل اوپن لاجک پروجیکٹ وچ شامل سارا مواد لوڈ کردی اے۔ ایہو جہے
ادارتی نوٹ، جے دکھائی دین، تاں ایہ دس دے نیں کہ فائل نوں ایس گل دا
خیال کیتے بغیر کمپائل کیتا گیا اے کہ مواد کِویں پیش ہونا چاہیدا اے۔
جے تسیں ایہ پڑھ سکدے او، تاں غالباً ایس پی ڈی ایف توں پڑھانا یا پڑھنا
\emph{مناسب نہیں}۔

اوپن لاجک پروجیکٹ کئی طریقے دیندا اے جنہاں نال پڑھاؤن یا اپنے آپ
پڑھن لئی زیادہ مناسب متن تیار کیتا جا سکدا اے۔ مثال لئی، پہلے توں طے
ترتیب وچ متن سارے منطقی عامل بنیادی ماندا اے تے ثبوتاں وچ ہر عامل
دیاں ساریاں صورتاں نوں حل کردا اے۔ پر ایہناں وچوں کجھ صورتاں نوں مشقاں
دے طور تے چھڈنا زیادہ چنگا اے۔ اوپن لاجک پروجیکٹ اُتے ہالے کم وی جاری
اے۔ سانجھے کم تے بہتری نوں ہلا شیری دین لئی مواد اوہدوں وی شامل کیتا
جاندا اے جدوں اوہ صرف مسودہ ہووے، اوہدیاں مشقاں نہ ہون، وغیرہ۔ کِسے
کورس لئی تیار کیتی پی ڈی ایف وچ ایہ حصے شامل نہیں ہون گے۔

پڑھاؤن تے پڑھائی لئی زیادہ مناسب پی ڈی ایف لبھن لئی
\href{http://builds.openlogicproject.org/}{اوپن لاجک دی ویب سائٹ اُتے دستیاب نمونے دے کورس}
ویکھو۔ اپنا متن تیار کرن لئی تسیں
\href{https://github.com/OpenLogicProject/OpenLogic/tree/master/courses/sample}{نمونے دی مرکزی بلڈ فائل}
توں شروع کر سکدے او، یا ہور نفیس تے اگلی سطح دیاں مثالاں لئی ایس متن
توں بنیاں درسی کتاباں دے ماخذ ویکھ سکدے او۔
\end{editorial}

\olimport[sets-functions-relations]{sets-functions-relations-complete}

\olimport[propositional-logic]{propositional-logic}

\olimport[first-order-logic]{first-order-logic}

\olimport[model-theory]{model-theory}

\olimport[computability]{computability}

\olimport[turing-machines]{turing-machines}

\olimport[incompleteness]{incompleteness}

\olimport[second-order-logic]{second-order-logic}

\olimport[lambda-calculus]{lambda-calculus}

\olimport[many-valued-logic]{many-valued-logic}

\olimport[normal-modal-logic]{normal-modal-logic}

\olimport[applied-modal-logic]{applied-modal-logic}

\olimport[intuitionistic-logic]{intuitionistic-logic}

\olimport[counterfactuals]{counterfactuals}

\olimport[set-theory]{set-theory}

\olimport[methods]{methods}

\olimport[history]{history}

\olimport[reference]{reference}

\end{document}

